\documentclass[11pt, a4paper]{article}

\usepackage[T1]{fontenc}
\usepackage[utf8]{inputenc}
\usepackage[spanish]{babel}
\usepackage[a4paper, margin=2.5cm]{geometry}
\usepackage{amsmath, amssymb}
\usepackage{booktabs}
\usepackage{enumitem}
\usepackage{listings}
\usepackage{xcolor}
\usepackage{hyperref}

\definecolor{codigobg}{RGB}{245,245,245}
\definecolor{keyword}{RGB}{0,100,150}
\lstset{language=R, basicstyle=\ttfamily\small, backgroundcolor=\color{codigobg},
  keywordstyle=\color{keyword}\bfseries, frame=single, breaklines=true, showstringspaces=false}

\setlength{\parindent}{0pt}
\pagestyle{plain}

\begin{document}

\begin{center}
{\large\bfseries PONTIFICIA UNIVERSIDAD JAVERIANA}\\[3pt]
Facultad de Ciencias Econ\'omicas y Administrativas\\[2pt]
Estad\'istica 2026-II --- Prof.\ Jaime Polanco Jim\'enez\\[6pt]
{\LARGE\bfseries Problem Set 5}\\[3pt]
{\large Pruebas $t$ de Student y Tama\~no del Efecto}\\[6pt]
Asignado: Septiembre 18 (Sesi\'on 5) \quad$|$\quad Entrega: Septiembre 25 (Sesi\'on 6)
\end{center}

\vspace{6pt}
\rule{\linewidth}{0.5pt}
\vspace{2pt}
{\centering\small Repositorio: \url{https://github.com/polanco-jaime/Curso-Estadistica-2026-3}\par}
\vspace{6pt}

\textbf{Instrucciones.} Este Problem Set es individual y se resuelve en R sobre
microdatos ICFES. La entrega es digital al inicio de la clase indicada.
El c\'odigo R debe incluirse.

\vspace{10pt}

\section*{Enunciado}

Los estudiantes de Negocios Internacionales (NI) deben alcanzar una competencia
s\'olida en ingl\'es para su desempe\~no profesional. En este Problem Set
evaluar\'a si el puntaje promedio de ingl\'es en Saber Pro difiere del umbral
de referencia nacional y si existen diferencias entre IES p\'ublicas y privadas.

\begin{enumerate}[label=\arabic*)]
  \item \textbf{Prueba $t$ de una muestra:} ¿el puntaje promedio de ingl\'es
        en Saber Pro (\texttt{MOD\_INGLES\_PUNT}) de los estudiantes de NI
        difiere significativamente de 160 puntos?
        \begin{enumerate}[label=\alph*)]
          \item Plantee las hip\'otesis nula y alternativa:
                \[
                  H_0: \mu = 160 \quad \text{vs.} \quad H_1: \mu \neq 160
                \]
          \item Aplique la prueba $t$ de una muestra usando \texttt{t.test()}
                en R.
          \item Reporte el estad\'istico $t$, los grados de libertad, el
                $p$-valor y el intervalo de confianza al 95\%.
          \item ¿Rechaza $H_0$ al nivel de significancia $\alpha = 0.05$?
                Interprete el resultado en contexto.
        \end{enumerate}

  \item \textbf{Prueba $t$ de dos muestras independientes:} ¿existe diferencia
        significativa en el puntaje de ingl\'es entre estudiantes de NI de
        instituciones p\'ublicas y privadas?
        \begin{enumerate}[label=\alph*)]
          \item Plantee $H_0$ y $H_1$ para la comparaci\'on de medias.
          \item Aplique la prueba $t$ de dos muestras (\texttt{t.test()})
                asumiendo varianzas desiguales (prueba de Welch).
          \item Reporte el estad\'istico $t$, los grados de libertad, el
                $p$-valor y la diferencia de medias.
          \item ¿Rechaza $H_0$ al 5\%? ¿Cu\'al grupo tiene mayor puntaje
                promedio?
        \end{enumerate}

  \item \textbf{Tama\~no del efecto ($d$ de Cohen):} calcule la $d$ de Cohen
        para cuantificar la magnitud de la diferencia p\'ublico--privado
        encontrada en el punto anterior.
        \begin{enumerate}[label=\alph*)]
          \item Use el paquete \texttt{effsize} o calcule manualmente:
                \[
                  d = \frac{\bar{X}_{\text{privada}} - \bar{X}_{\text{p\'ublica}}}
                           {s_{\text{pooled}}}
                \]
          \item Interprete el valor de $d$ seg\'un la convenci\'on de Cohen
                (peque\~no: $|d| \approx 0.2$; mediano: $|d| \approx 0.5$;
                grande: $|d| \approx 0.8$).
          \item ¿El efecto encontrado tiene relevancia pr\'actica adem\'as de
                significancia estad\'istica?
        \end{enumerate}

  \item \textbf{Curva de potencia:} eval\'ue c\'omo la potencia estad\'istica
        de la prueba $t$ de dos muestras depende del tama\~no de muestra.
        \begin{enumerate}[label=\alph*)]
          \item Usando la funci\'on \texttt{power.t.test()}, calcule la
                potencia de la prueba para detectar la diferencia observada
                (use $d$ de Cohen del punto anterior) con tama\~nos de muestra
                $n \in \{30, 50, 100, 200, 500\}$ por grupo.
          \item Grafique la curva de potencia $(n, \text{potencia})$.
          \item ¿Cu\'al es el tama\~no de muestra m\'inimo por grupo necesario
                para alcanzar una potencia de 0.80?
        \end{enumerate}

  \item \textbf{Conclusiones:} sintetice los hallazgos principales:
        \begin{enumerate}[label=\alph*)]
          \item ¿El puntaje promedio de NI difiere del umbral de 160 puntos?
          \item ¿Existe brecha p\'ublico--privado en ingl\'es? ¿Es grande?
          \item ¿El tama\~no de muestra actual es adecuado para detectar
                diferencias relevantes?
        \end{enumerate}
\end{enumerate}

\vspace{10pt}
\rule{\linewidth}{0.4pt}
\begin{center}
{\small Cada entrega completa suma $+0.0625$ a la nota final.}
\end{center}

\end{document}
